\documentclass[10pt]{article}
\usepackage[UTF8]{ctex}
\usepackage{amsmath, amssymb}
\usepackage{geometry}
\geometry{a4paper, margin=1.5cm}
\usepackage{enumitem}

\title{生活情景与数学建模原理}
\author{}
\date{}

\begin{document}
\maketitle

\section{成绩波动——为何中等生的分数起伏最大？}

每次大考后，班里总会出现两类学生：一类排名几乎纹丝不动，另一类则名次忽上忽下。一个反复被观察到的现象是，处于班级中游的同学，成绩波动往往最为剧烈；而顶尖学生和极低分学生的波动幅度反而较小。

为了解释这一现象，我们假设试卷可拆分为大量独立的得分机会（选择、填空、解答题步骤等），每个得分点仅存在“得满分”或“不得分”两种结果，服从伯努利分布。同时假设试卷由简单、中档、难题三类构成，权重分别为 $w_1=0.2,\ w_2=0.6,\ w_3=0.2$。该分布符合多数高中期中/期末考试的命题结构（中档题占主体以保证区分度，两端控制难度梯度）。设学生在三类题上的平均得分率为 $p_1,p_2,p_3$。中等生 $p_2\in[0.6,0.8],\ p_3\in[0.2,0.4]$；学霸 $p_i\to 1$；学困生 $p_i\to 0$。该区间划分基于教育测量学中的“难度-区分度”经验曲线。

基于中心极限定理，大量独立伯努利试验的加权和近似服从正态分布。总得分率 $R$ 的期望与方差可表示为：
$$
E(R) = w_1 p_1 + w_2 p_2 + w_3 p_3
$$
$$
D(R) \approx w_1 p_1(1-p_1) + w_2 p_2(1-p_2) + w_3 p_3(1-p_3)
$$
方差结构由 $f(p)=p(1-p)$ 主导。该二次函数在 $p=0.5$ 处取最大值 $0.25$，在 $p=0,1$ 处趋近于 $0$。因此，得分率越靠近 $0.5$ 的题型，对总波动的贡献越大。中等生因在中档题与难题上的 $p$ 值集中于 $0.2\sim0.8$ 区间，其 $D(R)$ 天然高于两极学生。

\begin{quote}
\textbf{结论：} 从统计结构上解释了“中等生波动最大”的必然性。
\end{quote}

\section{估分偏差——为何学霸估分更精准？}

考后对答案估分，有人和实际得分几乎一致，有人却常常相差十几分甚至几十分。一个稳定的课堂观察是：成绩越靠前，估分越“准”。这并不只是“自我认知”更强，也与阅卷中的两类随机因素有关：一类是“会做却丢分”——步骤不严谨、书写不规范、计算疏忽导致“自以为稳拿”的分点被扣；另一类是“不会却蹭分”——思路不完整但写对了关键句、列式有价值、结论碰巧对等，使“自以为全丢”的地方反而得到部分分。不同水平的学生，这两类事件发生的概率并不相同，于是估分偏差呈现出系统性规律。

将数学卷面满分 $150$ 分映射为得分率 $Y=S/150\in[0,1]$。学生自评得分率记为 $r=\widehat S/150$。设自认正确的点实际得分的概率（兑现率）为 $\alpha(r)=0.9+0.1r$，自认错误的点实际得分的概率（挽回率）为 $\beta(r)=0.1-0.1r$。参数依据：水平越高，书写规范、步骤严谨，自认正确的点失分概率越低（$r=1$ 时 $\alpha=1.0$）；同时高分段元认知准确，自认失分的点通常确实未得分，而中等生易对步骤分判断保守，故 $\beta(r)$ 随 $r$ 增大而减小（$r=0$ 时 $\beta=0.1$，$r=1$ 时 $\beta=0$）。

实际得分率 $Y$ 为两部分独立随机贡献之和。由期望线性性与方差独立可加性：
$$
E(Y) = r\alpha(r) + (1-r)\beta(r) = 0.2r^2 + 0.7r + 0.1
$$
系统偏差 $b(r)=E(Y)-r = 0.2r^2 - 0.3r + 0.1 = 0.1(2r-1)(r-1)$。$b(r)$ 是开口向上的二次函数，零点为 $r=0.5$ 与 $r=1.0$，顶点横坐标为 $r=0.75$。这意味着：$r\in(0,0.5)$ 时 $b(r)>0$，该分段学生整体倾向高估；$r\in(0.5,1)$ 时 $b(r)<0$，该分段学生整体倾向低估，并且 $r=0.75$ 时分数低估最为明显；$r=1.0$ 时 $b(r)=0$，顶尖生估分与实际期望完全对齐。

方差（发挥稳定性）为：
$$
D(Y)=r\alpha(1-\alpha) + (1-r)\beta(1-\beta)=0.09-0.08r-0.01r^2.
$$
函数 $D(Y)$ 在 $r\in[0,1]$ 上严格单调递减。这表示：水平越高，得分率的随机起伏越小，发挥越稳定。

均方误差（MSE）综合度量估分精度：$\text{MSE}(r)=D(Y)+[b(r)]^2$，同时惩罚系统偏差与随机波动，值越小估分越准。

\begin{quote}
\textbf{结论：} 学霸估分更精准的原因在于系统偏差更小（$r=1$ 时 $b=0$）且随机波动更小（方差单调递减）。
\end{quote}

\section{学习如逆水行舟——如何科学分配各科时间？}

勤奋的同学常面临时间分配困境：集中攻克弱科还是均衡分配？经验表明，从 $60$ 分提到 $70$ 分较易，从 $90$ 分提到 $95$ 分极难（边际递减）；同时，不学则退，数学衰退快，语文衰退慢（逆水行舟）。

假设每周投入固定时间 $t$ 后，科目分数趋于稳定状态。每周新增分数 = 学习效率 $\times$ 剩余空间 $\times$ 时间；每周衰退分数 = 衰退系数 $\times$ (当前分-基础分)。稳态时二者相等。为突出“策略比较”而非“精确预测”，采用情景化的示例参数，参数量级参考年级成绩跟踪与教师经验，保证各学科之间的相对大小符合常识。各科参数如下（满分语数英 150，选科 100）：

\begin{center}
\begin{tabular}{|c|c|c|c|c|c|}
\hline
科目 & $S_{\text{cur}}$（当前水平） & $M$（长期上限） & $S_0$（基础分） & $\eta$（学习效率） & $\theta$（衰退强度） \\
\hline
语文 & 105 & 120 & 95 & 0.01 & 0.05 \\
\hline
数学 & 90 & 120 & 70 & 0.05 & 0.20 \\
\hline
英语 & 105 & 125 & 95 & 0.02 & 0.10 \\
\hline
历史 & 80 & 90 & 70 & 0.03 & 0.08 \\
\hline
政治 & 75 & 85 & 70 & 0.03 & 0.10 \\
\hline
地理 & 75 & 85 & 70 & 0.03 & 0.10 \\
\hline
\end{tabular}
\end{center}

分类型解释：语文（语言积累型）的 $S_0$ 较高，$\eta$ 较小，$\theta$ 较小；数学（技能训练型）的 $S_0$ 较低，$\eta$ 较大，$\theta$ 较大；历史/政治/地理（记忆理解型）的 $\eta,\theta$ 取中间值，$M,S_0$ 与选科满分 100 的结构相匹配。假设各科学习时间互不干扰，总时间约束为线性。

把“每周给某一科投入固定时间 $t$”视为长期策略。一周的进步量 $\Delta S_+=\eta t(M-S)$，退步量 $\Delta S_-=\theta(S-S_0)$。稳态时 $\eta t(M-S^*) = \theta(S^*-S_0)$，解得：
$$
S^*(t) = \frac{\eta t M + \theta S_0}{\eta t + \theta}
$$
该函数满足 $S^*(0)=S_0$，$\lim_{t\to\infty}S^*(t)=M$，且一阶导数单调递减，天然刻画边际递减与衰退对抗。

为让模型结论更贴近学习决策，把每周时间分配拆成两步。第一步，计算每科“保底不下滑”的最少时间：令 $S^*(t)=S_{\text{cur}}$，解得 $t_{\min}=\frac{\theta\,(S_{\text{cur}}-S_0)}{\eta\,(M-S_{\text{cur}})}$。将所有 $t_{\min}$ 求和得 $T_{\min}$，则可用于“突破”的剩余时间为 $\Delta T=T-T_{\min}$。第二步，把 $\Delta T$ 全部加到某一科，即可计算突破后的稳态分数 $S^*(t_{\min}+\Delta T)$。若选择两科同时突破，最优分配满足等边际收益条件：设 $t_A=t_{A,\min}+x$，$t_B=t_{B,\min}+(\Delta T-x)$，则最优点满足 $\left.\frac{dS_A^*}{dt}\right|_{t=t_A}=\left.\frac{dS_B^*}{dt}\right|_{t=t_B}$，其中 $\frac{dS^*}{dt}=\frac{\eta\theta(M-S_0)}{\left(\eta t+\theta\right)^2}$。

\begin{quote}
\textbf{结论：} 先分配时间守住各科当前水平（避免下滑），再将剩余时间按“边际提分相等”的原则分配到最有潜力的科目上，可实现总分的最大化提升。
\end{quote}

\section{劳逸结合的科学依据——为何运动与思考不是浪费时间？}

连续刷题两小时后大脑反应变慢、错题率上升；起身运动二十分钟再回来，思路反而更清晰。“磨刀不误砍柴工”是否具有数学必然性？

假设每日课堂固定贡献 4 小时（不参与分配）；下午放学后有 1 小时可支配（唯一运动窗口）；晚自习有 3 小时可自由支配（仅能刷题或思考）。故 $t_1+t_2\le 3,\ 0\le t_3\le 1$。假设学习产出取决于知识积累 $T$、思维锐度 $S$、精力水平 $E$ 的乘积 $P=T\cdot S\cdot E$，任一因子过低将产生“木桶效应”。取 $T=4+0.9t_1$（$0.9<1$ 体现夜间刷题边际递减与疲劳累积），$S=1+0.2t_2$（深度思考每小时提升思维锐度 $0.2$），$E=1+0.15t_3$（运动增益依据《British Journal of Sports Medicine》荟萃分析，单次中等强度运动可使执行功能与专注力提升 12\%~18\%）。于是：
$$
P(t_1,t_2,t_3) = (4+0.9t_1)(1+0.2t_2)(1+0.15t_3)
$$
约束：$t_1+t_2\le 3,\ 0\le t_3\le 1,\ t_1,t_2,t_3\ge 0$。乘法结构使资源分配不再遵循简单加和最优，而需通过边际增长率均衡实现协同放大。

\begin{quote}
\textbf{结论：} 适当运动（$t_3>0$）和深度思考（$t_2>0$）能通过乘积效应提升总产出，即使减少了纯刷题时间（$t_1$），也可能使 $P$ 更大。这为“劳逸结合”提供了数学上的必然性。
\end{quote}

\end{document}